% ---------------------------------------------------------------------------
% ---------------------------------------------------------------------------
% Junção de templates encontrados no Overleaf
% facilitando a vida do aluno de pós da UFABC
% Modelo LaTex para preparação do documento final de Dissertação de Mestrado
% ninguém do programa de pós da informação validou se presta
% ---------------------------------------------------------------------------
% ---------------------------------------------------------------------------

\documentclass[
	% -- opções da classe memoir --
	12pt,					% tamanho da fonte
	openright,				% capítulos começam em pág ímpar (insere página vazia caso preciso)
	twoside,					% para impressão em verso e anverso. Oposto a oneside
	a4paper,					% tamanho do papel. 
	% -- opções da classe abntex2 --
	%chapter=TITLE,			% títulos de capítulos convertidos em letras maiúsculas
	%section=TITLE,			% títulos de seções convertidos em letras maiúsculas
	%subsection=TITLE,		% títulos de subseções convertidos em letras maiúsculas
	%subsubsection=TITLE,	% títulos de subsubseções convertidos em letras maiúsculas
	% -- opções do pacote babel --
	english,					% idioma adicional para hifenização
	%french,					% idioma adicional para hifenização
	%spanish,				% idioma adicional para hifenização
	brazil					% o último idioma é o principal do documento
	]{abntex2}

% ---------------------
% Pacotes OBRIGATÓRIOS
% ---------------------
\usepackage{lmodern}				% Usa a fonte Latin Modern			
\usepackage[T1]{fontenc}			% Selecao de codigos de fonte.
\usepackage[utf8]{inputenc}		% Codificacao do documento (conversão automática dos acentos)
\usepackage{lastpage}			% Usado pela Ficha catalográfica
\usepackage{indentfirst}			% Indenta o primeiro parágrafo de cada seção.
\usepackage{color}				% Controle das cores
\usepackage{graphicx,graphicx}	% Inclusão de gráficos
\usepackage{epsfig,subfig}		% Inclusão de figuras
\usepackage{microtype} 			% Melhorias de justificação
% ---------------------

% ---------------------
% Pacotes ADICIONAIS
% ---------------------
\usepackage{lipsum}						% Geração de dummy text
\usepackage{amsmath,amssymb,mathrsfs}	% Comandos matemáticos avançados 
\usepackage{setspace}  					% Para permitir espaçamento simples, 1 1/2 e duplo
\usepackage{verbatim}					% Para poder usar o ambiente "comment"
\usepackage{tabularx} 					% Para poder ter tabelas com colunas de largura auto-ajustável
\usepackage{afterpage} 					% Para executar um comando depois do fim da página corrente
\usepackage{url} 						% Para formatar URLs (endereços da Web)
\usepackage{tikz} 						% Para desenhar diagramas e figuras
\usetikzlibrary{positioning,arrows,shapes}
% ---------------------

% ---------------------
% Pacotes de CITAÇÕES
% ---------------------
\usepackage[brazilian,hyperpageref]{backref}	% Paginas com as citações na bibl
% \usepackage[alf]{abntex2cite}				% Citações padrão ABNT (alfa)
\usepackage[num]{abntex2cite}				% Citações padrão ABNT (numericas)
% ---------------------

% Configurações de CITAÇÕES para abntex2
% --- 
% CONFIGURAÇÕES DE PACOTES
% --- 

% ---
% Configurações do pacote backref
% Usado sem a opção hyperpageref de backref
\renewcommand{\backrefpagesname}{Citado na(s) página(s):~}
% Texto padrão antes do número das páginas
\renewcommand{\backref}{}
% Define os textos da citação
\renewcommand*{\backrefalt}[4]{
	\ifcase #1 %
		Nenhuma citação no texto.%
	\or
		Citado na página #2.%
	\else
		Citado #1 vezes nas páginas #2.%
	\fi}%
% ---

% Inclusão de dados para CAPA e FOLHA DE ROSTO (título, autor, orientador, etc.)
% ---
% Informações de dados para CAPA e FOLHA DE ROSTO
% ---
\titulo{Desenvolvimento de um Veículo Autônomo de Corrida Conforme as Regulamentações da Categoria "F1TENTH"
}
\autor{Gustavo Del Rio Lima, Julio Cesar Pereira de França}
\local{Santo André - SP}
\data{Xxxx de 2025}
\orientador{Filipe Ieda Fazanaro}
\instituicao{%
  Universidade Federal do ABC -- UFABC
  \par
  Centro de Engenharia, Modelagem e Ciências Sociais Aplicadas
  \par
  Programa de Graduação em Engenharia da Instrumentção Automação e Robótica}
\tipotrabalho{Dissertação (Mestrado)}
% O preambulo deve conter o tipo do trabalho, o objetivo,
% o nome da instituição e a área de concentração
\preambulo{\textbf{Dissertação de Conclusão de curso} apresentada ao Programa de Graduação em Engenharia da Instrumentção Automação e Robótica, como parte dos requisitos necessários para a obtenção do Título de Engenharia da Instrumentção Automação e Robótica.}
% ---

% Inclui Configurações de aparência do PDF Final
%  Configurações de aparência do PDF final
% NÃO ALTERAR!!!

% alterando o aspecto da cor azul
\definecolor{blue}{RGB}{41,5,195}

% informações do PDF
\makeatletter
\hypersetup{
     	%pagebackref=true,
		pdftitle={\@title}, 
		pdfauthor={\@author},
    		pdfsubject={\imprimirpreambulo},
	    pdfcreator={LaTeX with abnTeX2},
		pdfkeywords={abnt}{latex}{abntex}{abntex2}{trabalho acadêmico}, 
		colorlinks=true,       		% false: boxed links; true: colored links
    		linkcolor=blue,          	% color of internal links
    		citecolor=blue,        		% color of links to bibliography
    		filecolor=magenta,      		% color of file links
		urlcolor=blue,
		bookmarksdepth=4
} 
\makeatother
% --- 

% O tamanho da identação do parágrafo é dado por:
\setlength{\parindent}{1.3cm}

% Controle do espaçamento entre um parágrafo e outro:
\setlength{\parskip}{0.2cm}  % tente também \onelineskip

% ---------------------
% Compila o indice
% ---------------------
\makeindex
% ---------------------

%%%%%%%%%%%%%%%%%%%%%%%%%%%
%%  INICIO DO DOCUMENTO  %%
%%%%%%%%%%%%%%%%%%%%%%%%%%%
\begin{document}

% Retira espaço extra obsoleto entre as frases.
\frenchspacing

% ----------------------------------------------------------
% ELEMENTOS PRÉ-TEXTUAIS (Capa, Resumo, Abstract, etc.)
% ----------------------------------------------------------
\pretextual

% Capa
% ---
% Impressão da Capa
% ---
  \begin{capa}%
    \begin{figure}[h!]%
        \centering%
        \includegraphics[scale=1.2]{figs/logo.png}%
      \end{figure}%
    \center
	\ABNTEXchapterfont\large{Universidade Federal do ABC \\ Centro de Engenharia, Modelagem e Ciências Sociais Aplicadas \\ Programa de Graduação em Engenharia da Instrumentção Automação e Robótica}
	%\vspace{1.5cm}

    \vfill
    \ABNTEXchapterfont\bfseries\LARGE\imprimirtitulo
    \vfill

	%\vfill
	\ABNTEXchapterfont\large\imprimirautor
	\vfill
%
	Número de Ordem : XXXX
	
    \large\imprimirlocal, \large\imprimirdata

    \vspace*{1cm}
  \end{capa}
% ---

% Folha de rosto (o * indica que haverá a ficha bibliográfica)
\imprimirfolhaderosto*

% Imprimir Ficha Catalografica
% ---
% Ficha Catalográfica
% ---
% Isto é um exemplo de Ficha Catalográfica, ou ``Dados internacionais de
% catalogação-na-publicação''. Você pode utilizar este modelo como referência. 
% Porém, talvez a biblioteca lhe fornece um PDF
% com a ficha catalográfica definitiva após a defesa do trabalho. Quando estiver
% com o documento, salve-o como PDF no diretório do seu projeto e substitua todo
% o conteúdo de implementação deste arquivo pelo comando abaixo:
%
% \begin{fichacatalografica}
%     \includepdf{fig_ficha_catalografica.pdf}
% \end{fichacatalografica}
\begin{fichacatalografica}
	\vspace*{\fill}					% Posição vertical
	\hrule							% Linha horizontal
	\begin{center}					% Minipage Centralizado
	\begin{minipage}[c]{12.5cm}		% Largura
	
	\imprimirautor
	
	\hspace{0.5cm} \imprimirtitulo  / \imprimirautor. --
	\imprimirlocal, \imprimirdata-
	
	\hspace{0.5cm} \pageref{LastPage} p. : il. (algumas color.) ; 30 cm.\\
	
	\hspace{0.5cm} \imprimirorientadorRotulo~\imprimirorientador\\
	
	\hspace{0.5cm}
	\parbox[t]{\textwidth}{\imprimirtipotrabalho~--~\imprimirinstituicao,
	\imprimirdata.}\\
	
	\hspace{0.5cm}
		1. Palavra-chave1.
		2. Palavra-chave2.
		I. Orientador.
		II. Universidade xxx.
		III. Faculdade de xxx.
		IV. Título\\ 			
	
	\hspace{8.75cm} CDU 02:141:005.7\\
	
	\end{minipage}
	\end{center}
	\hrule
\end{fichacatalografica}
% ---

% Inserir Folha de Aprovação
% ---
% Assinaturas
% ---
% Isto é um exemplo de Folha de aprovação, elemento obrigatório da NBR
% 14724/2011 (seção 4.2.1.3). Você pode utilizar este modelo até a aprovação
% do trabalho. Após isso, substitua todo o conteúdo deste arquivo por uma
% imagem da página assinada pela banca com o comando abaixo:
%
% \includepdf{folhadeaprovacao_final.pdf}
%
\begin{folhadeaprovacao}

  \begin{center}
    {\ABNTEXchapterfont\large\imprimirautor}

    \vspace*{\fill}\vspace*{\fill}
    \begin{center}
      \ABNTEXchapterfont\bfseries\Large\imprimirtitulo
    \end{center}
    \vspace*{\fill}
    
    \hspace{.45\textwidth}
    \begin{minipage}{.5\textwidth}
        \imprimirpreambulo
    \end{minipage}%
    \vspace*{\fill}
   \end{center}
        
 % Isso na versao final do trabalho!!!       
   Trabalho aprovado. \imprimirlocal, 01 de janeiro de 2014:

   \assinatura{\textbf{\imprimirorientador} \\ Orientador} 
   \assinatura{\textbf{\imprimircoorientador} \\ Co-Orientador} 
   \assinatura{\textbf{Professor} \\ Convidado 1}
   \assinatura{\textbf{Professor} \\ Convidado 2}
   \assinatura{\textbf{Professor} \\ Convidado 3}
      
   \begin{center}
    \vspace*{0.5cm}
    {\large\imprimirlocal}
    \par
    {\large\imprimirdata}
    \vspace*{1cm}
  \end{center}
  
\end{folhadeaprovacao}
% ---

% Dedicatória
%% ---
% Dedicatória
% ---
\begin{dedicatoria}
   \vspace*{\fill}
   \centering
   \noindent
   \textit{ Aos verme que roeu as frias carnes de meu cadáver.} \vspace*{\fill}
\end{dedicatoria}
% ---

% Agradecimentos
% ---
% Agradecimentos
% ---
\begin{agradecimentos}



Agradeço a Xuxa, meus pais, cachorro, gato e papagaio, por ...

Agradeço ao meu orientador, XXXXXXXXX, por todos os conselhos, pela paciência e ajuda nesse período.

Aos meus amigos ...

Aos professores ...

À XXXXXX pelo apoio financeiro para realização deste trabalho de pesquisa.

\end{agradecimentos}
%% ---

% Epígrafe
% ---
% Epígrafe
% ---
\begin{epigrafe}
    \vspace*{\fill}
	\begin{flushright}
		\textit{``Não sei o que, \\
		          não sei o que,\\
                  não sei o que lá.''\\
		          (Autor Desconhecido)}
	\end{flushright}
\end{epigrafe}
% ---

% Resumo e Abstract
% ---
% RESUMOS
% ---

% RESUMO em português
\setlength{\absparsep}{18pt} % ajusta o espaçamento dos parágrafos do resumo
\begin{resumo}
 Segundo a ABNT, o resumo deve ressaltar o
 objetivo, o método, os resultados e as conclusões do documento. A ordem e a extensão
 destes itens dependem do tipo de resumo (informativo ou indicativo) e do
 tratamento que cada item recebe no documento original. O resumo deve ser
 precedido da referência do documento, com exceção do resumo inserido no
 próprio documento. Umas 10 linhas (\ldots) As palavras-chave devem figurar logo abaixo do
 resumo, antecedidas da expressão Palavras-chave:, separadas entre si por
 ponto e finalizadas também por ponto.

 \textbf{Palavras-chaves}: latex. abntex. editoração de texto.
\end{resumo}

% ABSTRACT in english
\begin{resumo}[Abstract]
 \begin{otherlanguage*}{english}
   This is the english abstract.

   \vspace{\onelineskip}
 
   \noindent 
   \textbf{Keywords}: latex. abntex. text editoration.
 \end{otherlanguage*}
\end{resumo}

% Lista de ilustrações
\pdfbookmark[0]{\listfigurename}{lof}
\listoffigures*
\cleardoublepage

% Lista de tabelas
\pdfbookmark[0]{\listtablename}{lot}
\listoftables*
\cleardoublepage

% Lista de abreviaturas e siglas
\begin{siglas}
	\item[ABNT] Associação Brasileira de Normas Técnicas
	\item[abnTeX] Normas para TeX
\end{siglas}

% Lista de símbolos
\begin{simbolos}
	\item[$ \Gamma $] Letra grega Gama
	\item[$ \Lambda $] Lambda
	\item[$ \zeta $] Letra grega minúscula zeta
	\item[$ \in $] Pertence
\end{simbolos}

% Inserir o SUMÁRIO
\pdfbookmark[0]{\contentsname}{toc}
\tableofcontents*
\cleardoublepage

% ----------------------------------------------------------
% ELEMENTOS TEXTUAIS (Capítulos)
% ----------------------------------------------------------
\textual
% Elementos textuais com numeração arábica
\pagenumbering{arabic}
% Reinicia a contagem do número de páginas
\setcounter{page}{1}

% Inclui cada capitulo da Dissertação
% ----------------------------------------------------------
% Introdução 
% Capítulo sem numeração, mas presente no Sumário
% ----------------------------------------------------------

\chapter*[Introdução]{Introdução}
\addcontentsline{toc}{chapter}{Introdução}

O desenvolvimento de veículos autônomos representa uma das fronteiras mais
promissoras da engenharia moderna, combinando avanços em robótica, inteligência
artificial, sistemas de controle e processamento de sinais
\cite{Karaman2011Optimal}. Neste contexto, a categoria F1TENTH emerge como uma
plataforma educacional valiosa, permitindo que pesquisadores trabalhem com os
mesmos princípios e desafios encontrados em veículos autônomos em escala real,
porém em um ambiente mais controlado e acessível \cite{OKelly2020F1TENTH}.

A relevância do desenvolvimento de veículos autônomos de corrida em escala
reduzida se estabelece em múltiplas dimensões. Do ponto de vista técnico, os
desafios enfrentados no desenvolvimento de um veículo F1TENTH espelham, em
menor escala, os problemas encontrados na indústria automotiva atual. Segundo
\cite{Babu2020Simulator}, a necessidade de desenvolver sistemas robustos de
percepção, planejamento de trajetória e controle em tempo real reflete
diretamente os desafios enfrentados por empresas que trabalham com veículos
autônomos em escala real.

A escolha da plataforma F1TENTH é particularmente apropriada por seu caráter
open-source e pela existência de uma comunidade internacional ativa. Como
destacam \cite{Agnihotri2020Teaching}, esta característica não apenas facilita
o acesso a recursos e conhecimentos, mas também permite que os resultados
obtidos possam ser compartilhados e utilizados por outros pesquisadores e
estudantes, multiplicando o impacto do trabalho desenvolvido. Do ponto de vista
prático, o desenvolvimento de um veículo autônomo em escala reduzida apresenta
vantagens significativas em termos de custos, segurança e facilidade de
experimentação. \cite{Wang2020LMPC} demonstram que a plataforma permite a
realização de testes e validações de forma mais ágil e segura do que seria
possível com veículos em escala real, possibilitando uma iteração mais rápida
no desenvolvimento e aperfeiçoamento dos algoritmos e sistemas.

Além disso, o projeto se alinha com as tendências atuais da indústria
automotiva, onde há uma crescente demanda por profissionais com conhecimento em
sistemas autônomos. \cite{Stachowicz2023FastRLAP} ressaltam que a experiência
adquirida no desenvolvimento deste tipo de trabalho pode contribuir
significativamente para a formação de profissionais melhor preparados para
atuar neste mercado em expansão.

Neste contexto, este trabalho tem como objetivo principal desenvolver um
veículo autônomo de corrida seguindo as especificações da categoria F1TENTH,
abordando aspectos desde a construção do modelo físico até a implementação dos
sistemas de percepção, planejamento de trajetória e controle. O desenvolvimento
será realizado considerando as questões de segurança e verificação formal
destacadas por \cite{Ivanov2020Verifying}, que enfatizam a importância da
validação rigorosa em sistemas autônomos.

O trabalho está organizado em capítulos que abordam: fundamentação teórica,
detalhando os conceitos básicos necessários; desenvolvimento do hardware,
descrevendo a plataforma mecânica e componentes eletrônicos; desenvolvimento do
software, apresentando a arquitetura e implementação dos algoritmos; e
resultados e discussão, analisando o desempenho do veículo e comparando com
trabalhos similares.

A estrutura apresentada permite uma abordagem sistemática do problema,
contribuindo para o avanço do conhecimento em uma área tecnológica estratégica,
enquanto promove a integração de diferentes áreas do conhecimento em
engenharia. O resultado esperado não é apenas um veículo autônomo funcional,
mas também um conjunto de conhecimentos e experiências que podem beneficiar
futuros projetos e pesquisas na área de sistemas autônomos.

\section*{Justificativa}\label{sec:Justificativa}
\addcontentsline{toc}{section}{Justificativa}

O desenvolvimento de veículos autônomos representa uma das fronteiras mais
promissoras da engenharia moderna, combinando avanços em robótica, inteligência
artificial, sistemas de controle e processamento de sinais. Neste contexto, a
categoria F1TENTH emerge como uma plataforma educacional valiosa, permitindo
que estudantes e pesquisadores trabalhem com os mesmos princípios e desafios
encontrados em veículos autônomos em escala real, porém em um ambiente mais
controlado e acessível.

A relevância deste trabalho se estabelece em múltiplas dimensões. Do ponto de
vista educacional, o desenvolvimento de um veículo autônomo de corrida oferece
uma oportunidade única de aplicar conhecimentos teóricos em um projeto prático
e desafiador. Os conceitos fundamentais de engenharia de controle,
processamento de sinais, programação em tempo real e robótica convergem em um
único projeto, proporcionando uma experiência de aprendizado integrada e
multidisciplinar.

No âmbito tecnológico, os desafios enfrentados no desenvolvimento de um veículo
F1TENTH espelham, em menor escala, os problemas encontrados na indústria
automotiva atual. A necessidade de desenvolver sistemas robustos de percepção,
planejamento de trajetória e controle em tempo real reflete diretamente os
desafios enfrentados por empresas que trabalham com veículos autônomos em
escala real.

Do ponto de vista prático, o desenvolvimento de um veículo autônomo em escala
reduzida apresenta vantagens significativas em termos de custos, segurança e
facilidade de experimentação. A plataforma permite a realização de testes e
validações de forma mais ágil e segura do que seria possível com veículos em
escala real, possibilitando uma iteração mais rápida no desenvolvimento e
aperfeiçoamento dos algoritmos e sistemas.

\section*{Motivação}\label{sec:motivacao}
\addcontentsline{toc}{section}{Motivação}

A escolha da plataforma F1TENTH é particularmente apropriada por seu caráter
open-source e pela existência de uma comunidade internacional ativa. Esta
característica não apenas facilita o acesso a recursos e conhecimentos, mas
também permite que os resultados obtidos possam ser compartilhados e utilizados
por outros pesquisadores e estudantes, multiplicando o impacto do trabalho
desenvolvido.

Além disso, o projeto se alinha com as tendências atuais da indústria
automotiva, onde há uma crescente demanda por profissionais com conhecimento em
sistemas autônomos. A experiência adquirida no desenvolvimento deste trabalho
pode contribuir significativamente para a formação de profissionais melhor
preparados para atuar neste mercado em expansão.

A competitividade inerente à categoria F1TENTH também adiciona um elemento
motivador importante ao projeto. A necessidade de otimizar o desempenho do
veículo para competições estimula a busca por soluções inovadoras e eficientes,
promovendo um ambiente de aprendizado dinâmico e desafiador.

Por fim, este trabalho se justifica pela sua contribuição para o avanço do
conhecimento em uma área tecnológica estratégica, oferecendo oportunidades de
desenvolvimento técnico e profissional, ao mesmo tempo em que promove a
integração de diferentes áreas do conhecimento em engenharia. O resultado
esperado não é apenas um veículo autônomo funcional, mas também um conjunto de
conhecimentos e experiências que podem beneficiar futuros projetos e pesquisas
na área de sistemas autônomos.

\section*{Objetivos}\label{sec:objetivos}
\addcontentsline{toc}{section}{Objetivos}

\lipsum[36]


% PARTE - Define a divisão do documento em partes (Não é obrigatório)
\part{Preparação da pesquisa}
%\chapter{Estado da Arte}\label{cap:estArte}

\lipsum[34]

\section*{Trabalhos Relacionados a Isto}\label{sec:primTrab}
\addcontentsline{toc}{section}{Trabalhos Relacionados a Isto}

\lipsum[34-36]
\chapter{Materiais e Métodos}\label{cap:ferramentas}

\section{Fundamentação Teórica}

\subsection{Odometria: Estimativa de Pose por Sensoriamento de Movimento}

Em sistemas de robótica móvel, a capacidade de estimar a pose do robô (posição
e orientação) ao longo do tempo é fundamental para a navegação autônoma. Uma
das técnicas mais básicas e amplamente utilizadas para esta finalidade é a
odometria. A odometria, em sua essência, é um método que utiliza dados de
sensores de movimento do próprio robô para estimar a sua pose relativa a um
ponto de partida conhecido \cite{thrun2005}. Diferentemente de sistemas de
localização global (GPS), a odometria é um sistema de localização relativo, que
acumula erros ao longo do tempo, mas que fornece uma estimativa contínua da
pose do robô em alta frequência, sendo crucial para diversas tarefas de
controle e planejamento em robótica móvel.

\subsubsection{Princípios da Odometria de Rodas}

Para robôs móveis com rodas, a forma mais comum de odometria é a odometria de
rodas. Esta técnica baseia-se na medição do movimento rotacional das rodas do
robô, tipicamente através de encoders rotativos acoplados aos eixos das rodas.
Os encoders fornecem pulsos elétricos a cada incremento de rotação da roda,
permitindo quantificar o deslocamento angular de cada roda.

Com base no conhecimento da geometria do robô (distância entre as rodas, raio
das rodas) e nas leituras dos encoders, é possível estimar o movimento do robô.
Através de modelos cinemáticos simples, o movimento rotacional das rodas é
convertido em deslocamento linear e angular do robô. Por exemplo, em um robô
diferencial, a diferença na rotação das rodas direita e esquerda indica a
rotação do robô, enquanto a média da rotação indica o movimento linear para
frente ou para trás.

\subsubsection{Odometria e Sistemas de Coordenadas no ROS2}

No contexto do ROS2, a odometria desempenha um papel fundamental na manutenção
do sistema de coordenadas \texttt{odom}. Conforme mencionado anteriormente, o
frame \texttt{odom} é um sistema de coordenadas local que fornece estimativas
de posição relativa. A odometria é tipicamente o principal componente
responsável por fornecer a transformação entre o frame \texttt{odom} e o frame
\texttt{base\_link}, representando a pose do robô em relação ao referencial de
odometria.

É importante notar que o frame \texttt{odom} não é um sistema de coordenadas global e fixo como o frame \texttt{map}. O frame \texttt{odom} é um referencial local que se move e deriva ao longo do tempo, acompanhando a estimativa de pose da odometria.

\subsubsection{Limitações da Odometria}

Apesar de sua importância, a odometria é inerentemente propensa a erros. Como a
pose é estimada através da integração de pequenos movimentos incrementais, os
erros de medição acumulam-se ao longo do tempo, resultando no fenômeno
conhecido como deriva (\textit{drift}). As principais fontes de erro na
odometria de rodas incluem:

\begin{itemize}
    \item \textbf{Deslizamento das rodas (wheel slippage):} Ocorre quando as rodas patinam ou deslizam sem rotacionar de forma precisa sobre a superfície, levando a leituras incorretas dos encoders.
    \item \textbf{Irregularidades no terreno:} Superfícies irregulares ou obstáculos podem causar movimentos não intencionais das rodas, gerando erros na estimativa de movimento.
    \item \textbf{Imperfeições mecânicas e calibração:} Variações no diâmetro das rodas, desalinhamento dos encoders ou erros na calibração do sistema podem introduzir erros sistemáticos na odometria.
    \item \textbf{Resolução dos encoders:} Encoders com baixa resolução podem não capturar pequenos movimentos rotacionais, limitando a precisão da odometria.
\end{itemize}

Devido à acumulação de erros, a odometria não é adequada para localização
global precisa a longo prazo. A pose estimada pela odometria torna-se cada vez
mais incerta com o tempo e a distância percorrida.

\subsubsection{Importância da Odometria na Navegação Robótica}

Apesar de suas limitações, a odometria continua sendo um componente essencial
em sistemas de navegação robótica. Suas principais vantagens incluem:

\begin{itemize}
    \item \textbf{Alta frequência de atualização:} A odometria fornece estimativas de pose em alta taxa, crucial para sistemas de controle em tempo real e planejamento reativo.
    \item \textbf{Independência de sensores externos:} A odometria utiliza apenas sensores internos do robô, funcionando em ambientes sem infraestrutura externa (como GPS ou beacons).
    \item \textbf{Baixo custo computacional:} Os algoritmos de odometria são computacionalmente leves, permitindo sua execução em tempo real em plataformas embarcadas.
\end{itemize}

Em sistemas de navegação autônoma robustos, a odometria é frequentemente
utilizada em conjunto com outras técnicas de sensoriamento e localização, como
SLAM (Simultaneous Localization and Mapping) e localização baseada em sensores
externos (e.g., LiDAR, câmeras). A odometria fornece uma estimativa inicial e
contínua da pose, que pode ser refinada e corrigida por algoritmos de
localização mais precisos, como o AMCL (Adaptive Monte Carlo Localization)
mencionado anteriormente, que frequentemente utiliza a odometria como uma de
suas fontes de informação.

\subsubsection{Odometria no ROS2 e Navigation2}

No ecossistema ROS2 e na biblioteca Navigation2, a odometria é um dado de
entrada fundamental para diversos módulos. O Nav2 espera receber dados de
odometria através de tópicos ROS, tipicamente no formato de mensagens
\texttt{nav\_msgs/Odometry}. Estes dados são utilizados por algoritmos de
localização, como AMCL, para estimar a pose do robô no mapa. Além disso, a
odometria pode ser combinada com dados de outros sensores, como IMUs (Unidades
de Medida Inercial), através de pacotes como o \texttt{robot\_localization},
para obter uma estimativa de pose mais robusta e precisa através de técnicas de
fusão de sensores.

\subsection{Base\_link: Frame de Referência Central do Robô}

Em arquiteturas robóticas, especialmente no contexto do ROS2 e da biblioteca
Navigation2, o frame \texttt{base\_link} assume um papel central como o sistema
de coordenadas fundamental atrelado ao corpo do robô. Conforme mencionado na
seção sobre sistemas de coordenadas \cite{Babu2020Simulator}, o
\texttt{base\_link} é um frame rígido fixado ao robô, servindo como ponto de
referência para a localização de todos os outros componentes e sensores do
sistema robótico.

\subsubsection{Definição e Significado do Base\_link}

O \texttt{base\_link} é definido como um sistema de coordenadas local e fixo em
relação ao corpo principal do robô. Tipicamente, a origem do
\texttt{base\_link} é localizada no centro geométrico do robô, ou em um ponto
de referência significativo para o design do robô. A orientação do
\texttt{base\_link} também é definida em relação ao robô, com os eixos X, Y e Z
alinhados com as direções de movimento e orientação do robô (por exemplo, X
para frente, Y para a esquerda, Z para cima).

É crucial entender que o \texttt{base\_link} é um frame abstrato e virtual. Ele não corresponde necessariamente a um componente físico específico do robô. Em vez disso, ele representa um ponto de referência conceitual que facilita a organização e a transformação de dados sensoriais e de controle dentro do sistema robótico.

\subsubsection{Importância do Base\_link na Arquitetura Robótica}

O \texttt{base\_link} é essencial por diversas razões na arquitetura de um robô
móvel. Como ponto de referência unificado, ele serve como um sistema de
coordenadas comum para integrar dados de diferentes sensores (encoders, IMUs,
LiDARs, câmeras) e atuadores do robô (motores, direção). As transformações de
todos os sensores e atuadores são tipicamente expressas em relação ao
\texttt{base\_link}, facilitando a fusão de dados e o controle coordenado do
robô.

Na representação da pose do robô, a posição no mundo (em frames como
\texttt{map} ou \texttt{odom}) é geralmente expressa como uma transformação em
relação ao \texttt{base\_link}. Isso significa que, ao se referir à "pose do
robô", geralmente se está falando da pose do \texttt{base\_link} em relação a
outro frame de referência.

Os módulos de navegação, como o Navigation2 no ROS2, utilizam o
\texttt{base\_link} como a representação da pose do robô para planejamento de
trajetória, controle e localização. Os comandos de velocidade e as informações
de sensores são interpretados e gerados em relação ao \texttt{base\_link}.

Em ambientes de simulação e ferramentas de modelagem CAD, o \texttt{base\_link}
simplifica a representação do robô como uma entidade coesa, facilitando a
definição de suas propriedades físicas e cinemáticas.

\subsubsection{Relação do Base\_link com Map e Odom}

O \texttt{base\_link} se relaciona diretamente com outros frames de coordenadas
cruciais na navegação robótica, como \texttt{map} e \texttt{odom}. A relação
entre estes frames segue a hierarquia \texttt{map} → \texttt{odom} →
\texttt{base\_link}. A transformação \texttt{map} → \texttt{odom} descreve a
pose do frame \texttt{odom} em relação ao frame global \texttt{map}, sendo
tipicamente estimada por algoritmos de localização global. A transformação
\texttt{odom} → \texttt{base\_link} descreve a pose do \texttt{base\_link} em
relação ao frame \texttt{odom}, sendo primariamente estimada pela odometria.

Em essência, o \texttt{base\_link} é o frame mais próximo do robô físico,
enquanto o \texttt{odom} fornece uma estimativa de pose relativa baseada em
sensores de movimento, e o \texttt{map} representa o sistema de coordenadas
global e fixo do ambiente. A transformação entre \texttt{map} e \texttt{odom}
permite corrigir os erros acumulados na odometria, ancorando a pose do robô no
mapa global.

\subsubsection{Base\_footprint: Projeção no Plano do Chão}

É importante mencionar também o frame \texttt{base\_footprint}, que é a projeção do \texttt{base\_link} no plano do chão. O \texttt{base\_footprint} é frequentemente utilizado em planejadores de trajetória e controladores de movimento que operam em um plano 2D. A transformação entre \texttt{base\_link} e \texttt{base\_footprint} é tipicamente uma translação vertical ao longo do eixo Z, representando a altura do \texttt{base\_link} em relação ao solo.

\subsubsection{Considerações Práticas e Utilização no ROS2}

No ROS2, a configuração correta do \texttt{base\_link} e suas transformações é
fundamental para o funcionamento adequado da Navigation2 e de outros pacotes
robóticos. Tipicamente, o frame \texttt{base\_link} é definido no arquivo URDF
(Unified Robot Description Format) do robô, que descreve a estrutura física e
cinemática do robô. As transformações entre \texttt{base\_link} e outros frames
são gerenciadas pelo pacote tf2, que publica e mantém as relações entre os
sistemas de coordenadas em tempo real.

Ao desenvolver aplicações robóticas no ROS2, é essencial garantir que o
\texttt{base\_link} esteja corretamente definido e que as transformações para
outros frames, especialmente \texttt{odom} e \texttt{map}, sejam precisas e
consistentes. A configuração e calibração adequadas dos sensores de odometria e
dos algoritmos de localização são cruciais para obter uma estimativa de pose
robusta e confiável do \texttt{base\_link} no ambiente.

\subsection{Conceitos Básicos de Robótica Móvel e Navegação ROS2}

A navegação autônoma de robôs móveis representa um dos desafios fundamentais da
robótica moderna, exigindo a integração de múltiplos subsistemas para permitir
que um robô se mova de forma segura e eficiente em seu ambiente. No contexto do
ROS2 (Robot Operating System 2), a biblioteca Navigation2 (Nav2) fornece uma
arquitetura modular e extensível para implementação de navegação autônoma
\cite{OKelly2020F1TENTH}.

\subsubsection{Sistemas de Coordenadas e Transformações}

Um conceito fundamental para a navegação robótica é o sistema de coordenadas e
suas transformações. No ROS2, o pacote tf2 (Transform Frame 2) é responsável
por manter e gerenciar as relações entre diferentes sistemas de coordenadas do
robô. Segundo \cite{Babu2020Simulator}, os principais frames utilizados na
navegação incluem:

\begin{itemize}
    \item \texttt{map}: Sistema de coordenadas global fixo, representando o ambiente de navegação
    \item \texttt{odom}: Sistema de coordenadas local que fornece estimativas de posição relativa
    \item \texttt{base\_link}: Sistema de coordenadas fixo ao centro do robô
    \item \texttt{base\_footprint}: Projeção do \texttt{base\_link} no plano do chão
\end{itemize}

A relação entre estes frames é fundamental para a localização e navegação do
robô. O Nav2 utiliza estas transformações para converter medições de sensores e
comandos de controle entre diferentes sistemas de coordenadas, permitindo uma
navegação precisa e coerente.

\subsubsection{Mapeamento e Localização}

O processo de SLAM (Simultaneous Localization and Mapping) é essencial para a
navegação autônoma, permitindo que o robô construa um mapa do ambiente enquanto
simultaneamente se localiza nele. \cite{Wang2020LMPC} destacam que o Nav2
suporta diferentes algoritmos de SLAM, incluindo:

\begin{itemize}
    \item \textbf{Mapeamento}: Criação de representações do ambiente através de mapas de ocupação
    \item \textbf{Localização}: Estimativa da pose do robô utilizando algoritmos como AMCL (Adaptive Monte Carlo Localization)
    \item \textbf{Loop Closure}: Detecção e correção de inconsistências no mapa quando revisitando áreas conhecidas
\end{itemize}

\subsubsection{Planejamento de Trajetória}

O planejamento de trajetória é estruturado em duas camadas principais, conforme
descrito por \cite{Karaman2011Optimal}:

\begin{itemize}
    \item \textbf{Planejador Global}: Responsável por encontrar um caminho livre de obstáculos do ponto atual até o objetivo, utilizando algoritmos como:
          \begin{itemize}
              \item A* para busca em grafos
              \item RRT (Rapidly-exploring Random Trees) para espaços contínuos
              \item Navfn para navegação baseada em campos potenciais
          \end{itemize}

    \item \textbf{Planejador Local}: Responsável por gerar comandos de velocidade em tempo real, considerando:
          \begin{itemize}
              \item Desvio de obstáculos dinâmicos
              \item Restrições cinemáticas do robô
              \item Otimização local da trajetória
          \end{itemize}
\end{itemize}

\subsubsection{Controle e Execução}

O sistema de controle no Nav2 é responsável por converter os planos de
trajetória em comandos de movimento para os atuadores do robô.
\cite{Stachowicz2023FastRLAP} destacam os principais componentes:

\begin{itemize}
    \item \textbf{Controladores de Trajetória}: Implementam algoritmos como Pure Pursuit, DWA (Dynamic Window Approach) ou TEB (Timed Elastic Band)
    \item \textbf{Recovery Behaviors}: Estratégias para recuperação de situações de falha
    \item \textbf{Camada de Segurança}: Garante operação segura através de monitores de colisão e limites de velocidade
\end{itemize}

\subsubsection{Arquitetura de Navegação}

A arquitetura do Nav2 é construída sobre o conceito de Behavior Trees, que
segundo \cite{Ivanov2020Verifying}, oferece várias vantagens:

\begin{itemize}
    \item \textbf{Modularidade}: Facilita a adição e remoção de comportamentos
    \item \textbf{Reusabilidade}: Permite a composição de comportamentos complexos a partir de blocos básicos
    \item \textbf{Flexibilidade}: Suporta diferentes estratégias de navegação através da reconfiguração da árvore
\end{itemize}

A árvore de comportamentos coordena a execução dos diferentes módulos do
sistema, incluindo:
\begin{itemize}
    \item Planejamento de trajetória
    \item Controle de movimento
    \item Recuperação de falhas
    \item Monitoramento de estado
\end{itemize}

Esta arquitetura modular permite que desenvolvedores adaptem e estendam o
sistema para diferentes aplicações e plataformas robóticas, mantendo a robustez
e confiabilidade necessárias para navegação autônoma.

\subsection{Sistemas de Controle e Navegação Autônoma}

A navegação autônoma de veículos requer uma arquitetura de controle robusta que
integre diferentes níveis de abstração, desde o controle de baixo nível dos
atuadores até o planejamento estratégico de alto nível. De acordo com
\cite{Wang2020LMPC}, essa hierarquia de controle pode ser estruturada em três
camadas principais que trabalham em conjunto para garantir uma navegação segura
e eficiente. O campo das corridas autônomas, revisado por
\cite{Betz2022Survey}, apresenta desafios únicos nesta integração,
especialmente ao operar nos limites dinâmicos do veículo.

\subsubsection{Controle de Baixo Nível}

O controle de baixo nível é responsável pela execução direta dos comandos nos
atuadores do veículo. Para veículos com geometria Ackermann, como é o caso do
F1TENTH, \cite{Babu2020Simulator} descrevem que o sistema de controle precisa
gerenciar dois aspectos fundamentais:

\begin{itemize}
    \item \textbf{Controle de Direção}: Responsável por ajustar o ângulo das rodas dianteiras através do servomotor. O modelo cinemático da direção Ackermann pode ser expresso pela equação:
          \begin{equation}
              \tan(\delta) = \frac{L}{R}
          \end{equation}
          Onde $\delta$ é o ângulo de esterçamento, $L$ é a distância entre eixos e $R$ é o raio de curvatura instantâneo.

    \item \textbf{Controle de Velocidade}: Gerencia a velocidade longitudinal do veículo através do motor de tração. O modelo dinâmico simplificado pode ser representado por:
          \begin{equation}
              \dot{v} = \frac{F_t - F_r}{m}
          \end{equation}
          Onde $v$ é a velocidade, $F_t$ é a força de tração, $F_r$ representa as forças resistivas e $m$ é a massa do veículo.
\end{itemize}

\subsubsection{Controle Preditivo Baseado em Modelo (MPC)}

O controle preditivo baseado em modelo representa uma abordagem avançada para
navegação autônoma. \cite{Stachowicz2023FastRLAP} explicam que o MPC opera
resolvendo um problema de otimização em tempo real que considera:

\begin{itemize}
    \item \textbf{Modelo do Sistema}: Representa a dinâmica do veículo através de equações diferenciais:
          \begin{equation}
              \dot{x} = f(x,u)
          \end{equation}
          Onde $x$ é o estado do sistema e $u$ são as entradas de controle.

    \item \textbf{Função Objetivo}: Define os objetivos de controle, tipicamente incluindo:
          \begin{itemize}
              \item Minimização do erro de trajetória
              \item Suavidade dos comandos de controle
              \item Eficiência energética
          \end{itemize}

    \item \textbf{Restrições}: Limitações físicas e operacionais:
          \begin{itemize}
              \item Limites de velocidade e aceleração
              \item Restrições de esterçamento
              \item Limites de segurança
          \end{itemize}
\end{itemize}

\subsubsection{Controle Adaptativo e Aprendizado}

O controle adaptativo tem se mostrado particularmente relevante para aplicações
de corrida autônoma. \cite{Karaman2011Optimal} destacam que sistemas
adaptativos podem ajustar seus parâmetros em tempo real para:

\begin{itemize}
    \item \textbf{Compensar Incertezas}: Adaptação a variações nas condições da pista e dinâmica do veículo através de:
          \begin{itemize}
              \item Estimação online de parâmetros
              \item Ajuste automático de ganhos de controle
              \item Compensação de perturbações (por exemplo, \cite{Jain2020Bayesrace} exploram o
                    uso de Processos Gaussianos para corrigir modelos cinemáticos simples com base
                    na experiência prévia)
          \end{itemize}

    \item \textbf{Otimização de Desempenho}: Melhoria contínua através de:
          \begin{itemize}
              \item Aprendizado por reforço iterativo
              \item Ajuste baseado em experiência prévia
              \item Adaptação a diferentes condições de corrida
          \end{itemize}
\end{itemize}

\subsubsection{Arquitetura de Segurança}

\cite{Ivanov2020Verifying} enfatizam a importância de uma arquitetura de segurança robusta que inclua:

\begin{itemize}
    \item \textbf{Monitoramento de Estado}: Sistema que avalia continuamente:
          \begin{itemize}
              \item Integridade dos sensores
              \item Estado dos atuadores
              \item Consistência das medições
          \end{itemize}

    \item \textbf{Sistema de Supervisão}: Responsável por:
          \begin{itemize}
              \item Detecção de falhas
              \item Ativação de comportamentos de recuperação
              \item Garantia de operação dentro de limites seguros
          \end{itemize}

    \item \textbf{Camada de Verificação}: Implementa:
          \begin{itemize}
              \item Verificação formal de propriedades de segurança
              \item Validação de comandos de controle
              \item Garantias de estabilidade
          \end{itemize}
\end{itemize}

\subsubsection{Integração dos Sistemas}

A integração efetiva destes diferentes níveis de controle é crucial para o
desempenho global do sistema. \cite{OKelly2020F1TENTH} propõem uma arquitetura
que coordena:

\begin{itemize}
    \item \textbf{Comunicação entre Camadas}: Através de:
          \begin{itemize}
              \item Interfaces bem definidas
              \item Protocolos de mensagem padronizados
              \item Sincronização temporal precisa
          \end{itemize}

    \item \textbf{Gerenciamento de Estados}: Incluindo:
          \begin{itemize}
              \item Transições suaves entre modos de operação
              \item Tratamento de condições de exceção
              \item Coordenação de subsistemas
          \end{itemize}

    \item \textbf{Otimização Global}: Considerando:
          \begin{itemize}
              \item Objetivos de alto nível
              \item Restrições de segurança
              \item Eficiência operacional
          \end{itemize}
\end{itemize}

Esta estrutura hierárquica de controle permite que o veículo autônomo navegue
de forma segura e eficiente, adaptando-se a diferentes condições e cenários de
corrida enquanto mantém garantias de segurança e desempenho.

\subsection{Sensoriamento e Percepção do Ambiente}

O sistema de percepção em veículos autônomos de corrida desempenha um papel
fundamental na compreensão do ambiente e na tomada de decisões em tempo real.
Este sistema precisa processar informações de diversos sensores para criar uma
representação precisa e atualizada do ambiente, permitindo uma navegação segura
em altas velocidades.

\subsubsection{Sensores Principais}

No contexto da plataforma F1TENTH, \cite{OKelly2020F1TENTH} descrevem que os
principais sensores utilizados podem ser categorizados em três grupos
fundamentais:

\begin{itemize}
    \item O \textbf{LiDAR} (Light Detection and Ranging) representa o sensor primário
          para percepção do ambiente. Este sensor emite pulsos laser e mede o tempo de
          retorno do sinal refletido, criando uma nuvem de pontos bidimensional que
          representa a geometria do ambiente ao redor do veículo. A medição do LiDAR pode
          ser expressa matematicamente como:
          \begin{equation}
              d_i = \frac{c \cdot \Delta t_i}{2}
          \end{equation}
          Onde $d_i$ é a distância medida para o ponto $i$, $c$ é a velocidade da luz, e $\Delta t_i$ é o tempo de voo do pulso laser.

    \item A \textbf{Unidade de Medição Inercial (IMU)} fornece dados cruciais sobre o
          movimento do veículo. Segundo \cite{Stachowicz2023FastRLAP}, a IMU integra
          acelerômetros e giroscópios para medir:
          \begin{equation}
              \begin{bmatrix}
                  a_x \\
                  a_y \\
                  a_z
              \end{bmatrix} =
              \begin{bmatrix}
                  \ddot{x} \\
                  \ddot{y} \\
                  \ddot{z}
              \end{bmatrix} +
              \begin{bmatrix}
                  g_x \\
                  g_y \\
                  g_z
              \end{bmatrix}
          \end{equation}
          Onde $a$ representa as acelerações medidas, $\ddot{x}$ as acelerações reais do veículo, e $g$ a aceleração gravitacional.

    \item Os \textbf{encoders} nas rodas fornecem odometria, medindo o deslocamento do
          veículo através da rotação das rodas. A relação entre pulsos do encoder e
          deslocamento linear pode ser expressa como:
          \begin{equation}
              \Delta s = \frac{2\pi r}{N} \cdot n
          \end{equation}
          Onde $r$ é o raio da roda, $N$ é o número de pulsos por revolução, e $n$ é o número de pulsos contados.
\end{itemize}

\subsubsection{Fusão Sensorial}

A fusão sensorial é essencial para combinar dados de diferentes sensores e
criar uma estimativa mais precisa do estado do veículo e do ambiente.
\cite{Babu2020Simulator} explicam que o Filtro de Kalman Estendido (EKF) é
comumente utilizado para este propósito, operando em duas etapas:

\begin{itemize}
    \item \textbf{Predição}:
          \begin{align}
              \hat{x}_{k|k-1} & = f(\hat{x}_{k-1|k-1}, u_k)   \\
              P_{k|k-1}       & = F_k P_{k-1|k-1} F_k^T + Q_k
          \end{align}

    \item \textbf{Atualização}:
          \begin{align}
              K_k           & = P_{k|k-1} H_k^T (H_k P_{k|k-1} H_k^T + R_k)^{-1} \\
              \hat{x}_{k|k} & = \hat{x}_{k|k-1} + K_k(z_k - h(\hat{x}_{k|k-1}))
          \end{align}
\end{itemize}

\subsubsection{Processamento de Dados LiDAR}

O processamento dos dados do LiDAR é particularmente importante para a
navegação em alta velocidade. \cite{Wang2020LMPC} descrevem as principais
etapas deste processamento:

\begin{itemize}
    \item \textbf{Filtragem de Ruído}: Remoção de medições espúrias através de filtros estatísticos e temporais.
    \item \textbf{Segmentação}: Identificação de diferentes elementos na nuvem de pontos, como paredes, obstáculos e espaços livres.
    \item \textbf{Extração de Características}: Detecção de elementos relevantes como:
          \begin{itemize}
              \item Bordas da pista
              \item Obstáculos dinâmicos
              \item Pontos de referência para localização
          \end{itemize}
\end{itemize}

\subsubsection{Localização e Mapeamento}

A localização precisa é crucial para navegação em alta velocidade.
\cite{Karaman2011Optimal} destacam que o sistema de localização deve integrar
múltiplas fontes de informação:

\begin{itemize}
    \item \textbf{Localização Baseada em Mapa}:
          \begin{equation}
              p(x_t | z_{1:t}, u_{1:t}, m) = \eta p(z_t | x_t, m) \int p(x_t | x_{t-1}, u_t) p(x_{t-1} | z_{1:t-1}, u_{1:t-1}, m) dx_{t-1}
          \end{equation}

    \item \textbf{Mapeamento Online}:
          \begin{equation}
              p(m | z_{1:t}, x_{1:t}) = \prod_i p(m_i | z_{1:t}, x_{1:t})
          \end{equation}
\end{itemize}

\subsubsection{Detecção e Rastreamento de Obstáculos}

Em ambientes de corrida, a detecção e rastreamento de outros veículos é
fundamental. \cite{Ivanov2020Verifying} descrevem que este processo envolve:

\begin{itemize}
    \item \textbf{Detecção}: Identificação de objetos dinâmicos através de:
          \begin{itemize}
              \item Segmentação de nuvem de pontos
              \item Classificação de objetos
              \item Estimação de posição e velocidade
          \end{itemize}

    \item \textbf{Rastreamento}: Acompanhamento de objetos ao longo do tempo usando filtros bayesianos:
          \begin{equation}
              p(x_k | z_{1:k}) = \frac{p(z_k | x_k) p(x_k | z_{1:k-1})}{p(z_k | z_{1:k-1})}
          \end{equation}
\end{itemize}

A integração eficiente destes sistemas de sensoriamento e percepção é crucial
para o desempenho global do veículo autônomo, permitindo uma navegação segura e
eficiente em ambientes dinâmicos de alta velocidade.

\subsection{Algoritmos de Planejamento de Trajetória}

O planejamento de trajetória em veículos autônomos de corrida representa um
desafio complexo que requer a consideração simultânea de múltiplos objetivos:
minimização do tempo de volta, garantia de segurança, respeito às restrições
dinâmicas do veículo e adaptação a condições variáveis da pista. Esta seção
explora os principais algoritmos e técnicas utilizados para enfrentar estes
desafios.

\subsubsection{Planejamento Global de Trajetória}

O planejamento global de trajetória estabelece um caminho ótimo considerando o
conhecimento completo da pista. Segundo \cite{Karaman2011Optimal}, este
processo pode ser formulado como um problema de otimização:

\begin{equation}
    \min_{x(t), u(t)} \int_0^T L(x(t), u(t))dt
\end{equation}

Sujeito às restrições:
\begin{align}
    \dot{x}(t)    & = f(x(t), u(t)) \\
    g(x(t), u(t)) & \leq 0
\end{align}

Onde $x(t)$ representa o estado do veículo, $u(t)$ os comandos de controle,
$L(\cdot)$ a função objetivo (tipicamente tempo ou energia), $f(\cdot)$ a
dinâmica do veículo, e $g(\cdot)$ as restrições do sistema.

\cite{Wang2020LMPC} destacam três abordagens principais para resolver este problema:

\begin{itemize}
    \item \textbf{Métodos Baseados em Amostragem (RRT*, PRM*)}: Estes algoritmos exploram o espaço de estados através de amostragem aleatória, construindo uma árvore ou grafo de trajetórias possíveis. O RRT* (Rapidly-exploring Random Tree*) garante otimalidade assintótica através de reconexão contínua da árvore:
          \begin{equation}
              d_{new} = \min_{x_{near}} \{cost(x_{init} \rightarrow x_{near}) + d(x_{near}, x_{new})\}
          \end{equation}

    \item \textbf{Métodos de Otimização Numérica}: Utilizam técnicas como programação quadrática sequencial (SQP) ou métodos de pontos interiores para encontrar trajetórias ótimas. A função objetivo típica inclui múltiplos termos:
          \begin{equation}
              J = w_t T + w_c \int_0^T \kappa(s)^2 ds + w_s \int_0^T \dot{s}(t)^2 dt
          \end{equation}
          Onde $T$ é o tempo total, $\kappa(s)$ é a curvatura da trajetória, e $\dot{s}(t)$ é a velocidade ao longo do caminho.

    \item \textbf{Métodos Baseados em Curvas Paramétricas}: Utilizam representações matemáticas como splines ou polinômios de Bernstein para gerar trajetórias suaves. Por exemplo, uma spline cúbica pode ser expressa como:
          \begin{equation}
              x(t) = a_0 + a_1t + a_2t^2 + a_3t^3
          \end{equation}
\end{itemize}

\subsubsection{Planejamento Local de Trajetória}

O planejamento local adapta a trajetória global em tempo real para lidar com
obstáculos dinâmicos e condições variáveis. \cite{Stachowicz2023FastRLAP}
apresentam três técnicas principais:

\begin{itemize}
    \item \textbf{Dynamic Window Approach (DWA)}: Considera um espaço de velocidades admissíveis $(v, \omega)$ baseado nas limitações dinâmicas do veículo:
          \begin{equation}
              V_a = \{(v,\omega) | v \in [v_{min}, v_{max}], \omega \in [\omega_{min}, \omega_{max}], \ddot{v} \leq a_{max}, \dot{\omega} \leq \alpha_{max}\}
          \end{equation}

    \item \textbf{Timed Elastic Band (TEB)}: Otimiza uma sequência de poses intermediárias considerando tempo e suavidade:
          \begin{equation}
              F_{total} = \sum_{i} (w_k F_{kin} + w_d F_{dyn} + w_t F_{time} + w_o F_{obs})
          \end{equation}

    \item \textbf{Model Predictive Contouring Control (MPCC)}: Combina seguimento de trajetória com otimização de tempo através de um horizonte deslizante:
          \begin{equation}
              \min_{u(\cdot)} \sum_{k=0}^{N-1} (l_{lag,k}^2 + l_{con,k}^2 - \rho v_{k})
          \end{equation}
          Outras abordagens, como as comparadas por \cite{Misra2024Report}, utilizam RRT* ou campos potenciais gaussianos para o desvio de obstáculos local.
\end{itemize}

\subsubsection{Otimização de Linha de Corrida}

\cite{OKelly2020F1TENTH} explicam que a otimização da linha de corrida envolve encontrar o caminho que minimiza o tempo de volta considerando as características dinâmicas do veículo. Este processo considera:

\begin{itemize}
    \item \textbf{Maximização da Velocidade em Curvas}: A velocidade máxima em uma curva é limitada pela força lateral máxima:
          \begin{equation}
              v_{max} = \sqrt{\mu g R}
          \end{equation}

    \item \textbf{Otimização de Pontos de Entrada e Saída}: A posição lateral na pista é otimizada para maximizar a velocidade média:
          \begin{equation}
              y(s) = \sum_{i=0}^N a_i \phi_i(s)
          \end{equation}
          Onde $\phi_i(s)$ são funções base e $a_i$ são coeficientes a serem otimizados.
\end{itemize}

\subsubsection{Considerações de Segurança e Robustez}

\cite{Ivanov2020Verifying} enfatizam a importância de garantir segurança e robustez no planejamento de trajetória através de:

\begin{itemize}
    \item \textbf{Margens de Segurança Adaptativas}:
          \begin{equation}
              d_{safe}(v) = d_{min} + k_v v + k_a a
          \end{equation}

    \item \textbf{Verificação de Colisão Contínua}:
          \begin{equation}
              d_{obs}(t) > d_{safe}(v(t)) \quad \forall t \in [0,T]
          \end{equation}

    \item \textbf{Consideração de Incertezas}:
          \begin{equation}
              P(collision) \leq \epsilon
          \end{equation}
          Técnicas específicas como o teste de estresse baseado em RRT podem ser usadas para encontrar cenários de falha em manobras complexas como ultrapassagens \cite{Bak2021Stress}.
\end{itemize}

A integração efetiva destes diferentes aspectos do planejamento de trajetória é
crucial para o desenvolvimento de um sistema de corrida autônomo capaz de
operar de forma segura e eficiente em condições de alta velocidade.

\section{Considerações Finais}

\lipsum[23]

\chapter{Construção do Veículo F1TENTH}

Este capítulo detalha o processo de desenvolvimento e construção do veículo
para a competição F1TENTH, abrangendo desde a seleção de componentes até a
integração dos sistemas e validação do protótipo. Cada decisão de projeto foi
tomada considerando tanto as regulamentações da competição quanto os objetivos
de desempenho estabelecidos.

O objetivo principal foi criar um veículo que não apenas atendesse às
especificações técnicas obrigatórias, mas que também oferecesse uma plataforma
flexível para o desenvolvimento e teste de algoritmos avançados de percepção,
planejamento e controle.

\section{Seleção de Componentes}

A seleção de componentes foi realizada considerando múltiplos fatores como
conformidade com as regulamentações, desempenho, disponibilidade, custo e
potencial de integração. Cada componente foi avaliado quanto à sua contribuição
para os objetivos gerais do projeto e complementaridade com outros elementos do
sistema.

\subsection{Chassi e Sistema Mecânico}

Para a estrutura base do veículo, foi selecionado o chassi modelo D5 da marca
Emida, específico para a escala 1:10. A escolha deste chassi foi baseada em
suas características e compatibilidade com os requisitos da competição F1TENTH.

As especificações técnicas principais do chassi D5 incluem:
\begin{itemize}
      \item Dimensões: Comprimento de 38.2 cm, largura de 19.5 cm e distância entre eixos
            de 257 mm.
      \item Peso: Aproximadamente 0.78 kg (sem eletrônica), proporcionando uma base leve.
      \item Materiais: Estrutura principal em fibra de carbono, com componentes mecânicos
            em metal, oferecendo um bom equilíbrio entre rigidez e peso.
      \item Configuração: Posição do motor central (mid-motor), favorecendo a distribuição
            de peso.
      \item Suspensão: Sistema independente na dianteira e duplo braço (double wishbone) na
            traseira, com amortecedores ajustáveis, permitindo adaptação a diferentes
            condições de pista.
      \item Compatibilidade: Projetado para superfícies on-road e compatível com
            componentes eletrônicos padrão da escala 1:10.
      \item Ajustabilidade: Múltiplas posições para montagem da bateria e do motor,
            permitindo otimizar o centro de gravidade.
\end{itemize}

A estrutura em fibra de carbono garante a rigidez necessária para suportar as
cargas dinâmicas durante a operação em alta velocidade, enquanto a configuração
mid-motor e as opções de ajuste contribuem para a otimização do balanceamento e
da dirigibilidade do veículo.

\subsection{Sistema de Propulsão}

Para o sistema de propulsão, foi selecionado o motor EZRUN 3665SD-3200KV-G3, um
motor brushless DC que atende às especificações da competição, oferecendo
excelente relação entre desempenho e conformidade regulatória.

\subsubsection{Motor EZRUN 3665SD-3200KV-G3}

O motor EZRUN 3665SD-3200KV-G3 foi escolhido por suas características técnicas
alinhadas com os requisitos da competição:

\begin{itemize}
      \item Rating KV de 3200KV, dentro do limite regulamentar equivalente ao Velineon
            3500KV
      \item Compatibilidade com baterias 2-3S LiPo, comum em veículos RC em escala 1:10
      \item Dimensões físicas (37mm de diâmetro e 65.8mm de comprimento) adequadas para
            instalação no chassi padrão
      \item Construção robusta com carcaça de alumínio usinado em CNC e rolamentos de alta
            precisão
      \item Sistema de refrigeração otimizado para operação contínua em alta carga
      \item Classificação IP-67, oferecendo proteção contra água e poeira, importante para
            a durabilidade durante testes e competições
\end{itemize}

O motor apresenta um equilíbrio ideal entre torque e velocidade, fundamental
para o desempenho tanto em acelerações rápidas quanto em trechos de velocidade
constante nos circuitos da competição. A compatibilidade com sensores hall
integrados também facilita o controle preciso de velocidade necessário para
algoritmos de navegação autônoma.

\subsection{Sistema de Controle Eletrônico}

Para o controle eletrônico do motor, foi escolhido o Makerbase VESC MINI 6.7,
um controlador avançado que permite implementação de estratégias sofisticadas
de controle do motor.

\subsubsection{Makerbase VESC MINI 6.7}

O VESC MINI 6.7 foi selecionado pelas seguintes características:

\begin{itemize}
      \item Dimensões compactas (46.5mm x 36.5mm x 12mm) e peso reduzido (aproximadamente
            30g), ideais para aplicações em veículos em escala 1:10
      \item Ampla faixa de tensão operacional (8V-60V DC), compatível com diversos sistemas
            de bateria
      \item Capacidade de corrente contínua de 40A, suficiente para o motor selecionado
      \item Microcontrolador STM32F405RGT6, oferecendo poder de processamento para
            implementação de algoritmos avançados de controle
      \item Suporte para operação com e sem sensores, permitindo flexibilidade no modo de
            controle
      \item Capacidade de frenagem regenerativa, importante para manobras precisas
      \item Múltiplas interfaces de comunicação (USB, CAN, UART, PPM, I2C), facilitando a
            integração com a unidade de processamento principal
      \item Proteções integradas contra sobretensão, subtensão, sobrecorrente e
            sobretemperatura
\end{itemize}

A escolha deste controlador permite a implementação de controle vetorial de
campo orientado (FOC), resultando em operação mais eficiente e precisa do
motor. A capacidade de programação através do VESC Tool possibilita ajustes
finos nos parâmetros de controle, adaptando o comportamento do veículo para
diferentes seções da pista e estratégias de corrida.

\subsection{Sistema de Sensoriamento}

O principal sensor para percepção do ambiente é o LiDAR YDLIDAR X4PRO, que
fornece dados de distância em 360 graus para navegação e detecção de
obstáculos.

\subsubsection{YDLIDAR X4PRO}

O LiDAR YDLIDAR X4PRO foi escolhido por suas características técnicas que
atendem aos requisitos da competição:

\begin{itemize}
      \item Varredura omnidirecional de 360°, permitindo detecção completa do ambiente ao
            redor do veículo
      \item Alcance de medição de 0.12-10m, adequado para ambientes de competição indoor
      \item Frequência de amostragem de até 5000Hz, fornecendo dados detalhados para
            algoritmos de percepção
      \item Frequência de varredura ajustável entre 6-12Hz, permitindo equilíbrio entre
            resolução e tempo de resposta
      \item Resolução angular entre 0.43° e 0.86° (dependendo da frequência), suficiente
            para detecção precisa dos limites da pista e outros veículos
      \item Alta resistência a interferências luminosas (até 40000 Lux), importante para
            ambientes com iluminação variável
      \item Classificação de segurança laser Class I, atendendo aos requisitos de segurança
            da competição
      \item Dimensões compactas e peso de 178g, facilitando a integração no veículo
\end{itemize}

Este sensor fornece dados cruciais para algoritmos de mapeamento, localização,
planejamento de trajetória e detecção de outros veículos durante as corridas
head-to-head. Sua capacidade de detecção omnidirecional é particularmente
importante para a implementação de estratégias defensivas e ofensivas durante
ultrapassagens.

\subsection{Unidade de Processamento}

Para processamento dos algoritmos de percepção, planejamento e controle, foi
selecionada a Raspberry Pi 4B como unidade computacional principal.

\subsubsection{Raspberry Pi 4B}

A Raspberry Pi 4B foi escolhida como plataforma de processamento principal
pelas seguintes razões:

\begin{itemize}
      \item Processador quad-core Cortex-A72 (ARM v8) 64-bit operando a 1.5GHz, oferecendo
            poder computacional suficiente para algoritmos de tempo real
      \item Disponibilidade de 4GB ou 8GB de RAM, dependendo da configuração, permitindo
            processamento de dados sensoriais complexos
      \item Múltiplas interfaces de comunicação (USB 3.0, Ethernet Gigabit, Wi-Fi,
            Bluetooth, GPIO), facilitando a integração com diversos sensores e atuadores
      \item Dimensões compactas e peso reduzido, adequados para instalação em veículos em
            escala 1:10
      \item Amplo suporte para ROS2 (Robot Operating System 2), framework utilizado para
            implementação dos algoritmos de navegação autônoma
      \item Comunidade ativa e extensa documentação, facilitando o desenvolvimento
      \item Custo acessível em comparação com outras plataformas de desempenho similar
\end{itemize}

A Raspberry Pi 4B executa o sistema operacional Ubuntu 22.04 LTS com ROS2
Humble Hawksbill, fornecendo um ambiente robusto para desenvolvimento e
execução dos algoritmos de controle. A integração com o VESC e o LiDAR é
realizada através de interfaces seriais e pacotes ROS2 específicos, criando uma
arquitetura de software modular e extensível.

\subsection{Sistema de Alimentação}

O sistema de alimentação do veículo foi projetado com duas fontes de energia
distintas para garantir a estabilidade operacional: uma bateria de alta
descarga para o sistema de propulsão e um power bank para a eletrônica de
controle e processamento.

\subsubsection{Bateria de Propulsão: HRB LiPo 3S 5000mAh 50C}

A alimentação principal do motor EZRUN e do controlador VESC é fornecida por
uma bateria HRB LiPo 3S de 5000mAh. Suas características são:
\begin{itemize}
      \item Tipo: Polímero de Lítio (LiPo)
      \item Configuração: 3 células em série (3S)
      \item Tensão Nominal: 11.1V
      \item Capacidade: 5000mAh
      \item Taxa de Descarga: 50C contínua (250A), 100C pico
      \item Dimensões: 155mm x 48mm x 24mm
      \item Peso: 376g
\end{itemize}
Esta bateria foi escolhida por sua alta capacidade e taxa de descarga, capaz de
fornecer a corrente necessária para o motor durante acelerações intensas,
garantindo o desempenho requerido na competição.

\subsubsection{Alimentação da Eletrônica: Essager Power Bank 15000mAh 65W}

Para alimentar a Raspberry Pi 4B, o LiDAR, servo motor e outros periféricos de
baixa tensão, foi utilizado um power bank Essager modelo Xuneng de 15000mAh.
Suas características relevantes incluem:
\begin{itemize}
      \item Capacidade: 15000mAh
      \item Potência Máxima de Saída: 65W
      \item Portas: 2x USB-C, 1x USB-A
      \item Peso: Aproximadamente 300g
      \item Funcionalidade Adicional: Display LED e cabo USB-C integrado.
\end{itemize}
O uso de um power bank dedicado para a eletrônica desacopla a alimentação da
Raspberry Pi das flutuações de tensão causadas pelo motor, garantindo uma
operação mais estável e prevenindo reinicializações inesperadas da unidade de
processamento. A capacidade de 15000mAh oferece autonomia suficiente para
sessões prolongadas de teste e operação.

\section{Projeto e Fabricação de Peças Customizadas}

Para acomodar os diversos componentes e garantir a integridade estrutural do
veículo, várias peças customizadas foram projetadas e fabricadas utilizando
tecnologia de impressão 3D.

\subsection{Suportes e Estruturas Principais}

[Esta seção será completada posteriormente com detalhes sobre o suporte para a unidade de processamento e controlador eletrônico.]

\subsection{Sistema de Montagem e Integração}

[Esta seção será completada posteriormente com informações sobre como as peças customizadas se integram ao chassi e entre si, incluindo considerações sobre manutenção e acesso aos componentes.]

\subsection{Processo de Fabricação}

[Esta seção será completada posteriormente com detalhes sobre os parâmetros de impressão 3D, materiais utilizados e técnicas de pós-processamento.]

\section{Arquitetura de Hardware do Veículo}

 [Esta seção será completada posteriormente com um diagrama detalhado da arquitetura de hardware e explicações sobre as conexões entre componentes.]

\section{Arquitetura de Software do Veículo}

A arquitetura de software do veículo é baseada em ROS2, implementando uma
estrutura modular que facilita o desenvolvimento, teste e otimização dos
algoritmos de controle.

\subsection{Integração do VESC com ROS2}

A integração do controlador VESC MINI 6.7 com o sistema ROS2 é realizada
utilizando os pacotes ROS2 mantidos pela comunidade F1TENTH, especificamente o
metapacote \texttt{vesc}\footnote{Repositório F1TENTH VESC ROS2 Driver:
      \url{https://github.com/f1tenth/vesc/tree/ros2}}. Este metapacote inclui:
\begin{itemize}
      \item \texttt{vesc\_driver}: O nó principal que gerencia a comunicação serial com o
            VESC via USB (tipicamente na porta \texttt{/dev/ttyACM0}).
      \item \texttt{vesc\_msgs}: Define os tipos de mensagens customizadas para
            comunicação com o VESC, como \texttt{VescStateStamped} para telemetria e
            comandos.
      \item \texttt{vesc\_ackermann}: Converte comandos de \texttt{AckermannDriveStamped} (velocidade
            linear e ângulo de direção) em comandos de baixo nível compreensíveis pelo
            VESC (velocidade do motor em ERPM e posição do servo).
\end{itemize}

A comunicação se dá da seguinte forma:
\begin{itemize}
      \item Publicação de Estado: O nó \texttt{vesc\_driver} lê dados do VESC via serial e
            publica a telemetria (tensão, corrente, ERPM, temperatura, etc.) no tópico
            \texttt{/sensors/core} (tipo \texttt{vesc\_msgs/VescStateStamped}) e a posição
            do servo (se configurado) em \texttt{/sensors/servo\_position\_command}.
      \item Recebimento de Comandos: O nó \texttt{vesc\_ackermann} recebe comandos de alto
            nível de velocidade e direção no tópico \texttt{/drive} (tipo
            \texttt{ackermann\_msgs/AckermannDriveStamped}).
      \item Conversão e Envio: \texttt{vesc\_ackermann} converte os comandos de Ackermann
            em comandos de ERPM para o motor e posição para o servo, publicando-os nos
            tópicos \texttt{/commands/motor/speed} e \texttt{/commands/servo/position}
            respectivamente.
      \item Atuação: O nó \texttt{vesc\_driver} subscreve a esses tópicos de comando e
            envia as instruções correspondentes ao VESC via comunicação serial.
\end{itemize}

A configuração do sistema, incluindo a porta serial, ganhos de conversão
(velocidade para ERPM, ângulo para posição do servo) e limites operacionais, é
gerenciada através de arquivos de parâmetros YAML (como o
\texttt{vesc\_config.yaml} no pacote \texttt{vesc\_driver}), carregados durante
a inicialização dos nós com \texttt{ros2 launch}.

Essa arquitetura modular permite que os algoritmos de planejamento e controle
operem em um nível de abstração mais alto, enviando comandos
\texttt{AckermannDriveStamped}, enquanto os pacotes \texttt{vesc} cuidam da
comunicação de baixo nível e da conversão necessária para o hardware
específico.

\subsection{Integração do LiDAR com ROS2}

A integração do sensor LiDAR YDLIDAR X4PRO com o sistema ROS2 é realizada
através do pacote oficial \texttt{ydlidar\_ros2\_driver}\footnote{Repositório
      oficial do driver: \url{https://github.com/YDLIDAR/ydlidar_ros2_driver}}. Este
driver atua como uma ponte entre o hardware do sensor e o framework ROS2,
permitindo que os dados de varredura laser sejam processados e utilizados pelos
algoritmos de navegação.

Principais aspectos da integração:
\begin{itemize}
      \item Comunicação: O driver estabelece comunicação com o LiDAR através da porta
            serial (tipicamente \texttt{/dev/ttyUSB0} ou um alias como
            \texttt{/dev/ydlidar}) na Raspberry Pi. Os parâmetros de comunicação, como
            \texttt{baudrate} (230400 para o X4PRO) e \texttt{port}, são configuráveis.
      \item Publicação de Dados: O driver processa os dados brutos do LiDAR e os publica no
            tópico ROS2 \texttt{/scan}. A mensagem publicada é do tipo
            \texttt{sensor\_msgs/LaserScan}, um formato padrão em ROS para dados de
            varredura laser 2D. Esta mensagem contém informações como ângulos, distâncias e
            intensidades (se aplicável) para cada ponto da varredura.
      \item Configuração: Parâmetros específicos do LiDAR, como faixa angular
            (\texttt{angle\_min}, \texttt{angle\_max}), limites de alcance
            (\texttt{range\_min}, \texttt{range\_max}), frequência de varredura
            (\texttt{frequency}), e \texttt{frame\_id} (identificador da coordenada TF,
            e.g., \texttt{laser\_frame}), são definidos em arquivos de configuração (YAML)
            ou diretamente em arquivos de inicialização (arquivos \texttt{.launch.py}).
      \item Execução: O nó do driver é iniciado utilizando \texttt{ros2 launch},
            tipicamente com um arquivo \texttt{launch.py} customizado que carrega os
            parâmetros corretos para o X4PRO (e.g., \texttt{ydlidar\_launch.py}
            referenciado no repositório).
      \item Serviços: O driver também oferece serviços ROS2 básicos, como
            \texttt{start\_scan} e \texttt{stop\_scan}, permitindo ligar e desligar o
            sensor programaticamente.
\end{itemize}

Essa integração fornece os dados de percepção essenciais, no formato
\texttt{sensor\_msgs/LaserScan}, para os nós ROS2 responsáveis pelo mapeamento
SLAM (Simultaneous Localization and Mapping), localização (como AMCL --
Adaptive Monte Carlo Localization), planejamento de trajetória e detecção de
obstáculos. A padronização da mensagem facilita a interoperabilidade com outros
pacotes ROS2 do ecossistema de navegação.

\subsection{Algoritmos de Percepção e Controle}

[Esta seção será completada posteriormente com informações sobre os algoritmos implementados para mapeamento, localização, planejamento de trajetória e controle do veículo.]

\chapter{Regulamentações e Dinâmica da Competição F1TENTH}

A categoria F1TENTH estabelece um framework competitivo rigoroso através de
regulamentações técnicas específicas que moldam diretamente o desenvolvimento
dos veículos e as estratégias de corrida~\cite{F1TENTHRules2023}.
Diferentemente de outras competições de robótica, a F1TENTH enfatiza o
desenvolvimento algorítmico em detrimento de vantagens puramente mecânicas,
criando um ambiente que aproxima os desafios enfrentados por veículos autônomos
em escala real~\cite{OBrienF1TENTH2020}.

\section{Especificações Técnicas Obrigatórias e suas Implicações}

\subsection{Plataforma mecânica e propulsão}

As regulamentações dividem a competição em duas categorias distintas: Classe
Restrita e Classe Aberta. Na Classe Restrita, foco principal da competição, o
chassi deve obrigatoriamente ser em escala 1:10, preferencialmente baseado em
modelos Traxxas, permitindo tanto configurações 4WD quanto 2WD
~\cite{F1TENTHRules2023}. Esta padronização não é arbitrária -- visa garantir
que todas as equipes enfrentem desafios físicos comparáveis, evitando vantagens
derivadas exclusivamente de plataformas superiores.

O sistema de propulsão também é rigidamente controlado, limitando-se a motores
DC sem escovas com especificações máximas equivalentes ao Velineon 3500 kV. A
escolha do motor EZRUN 3665SD-3200KV-G3 para o projeto representa uma
implementação adequada, pois seu rating de 3200KV está em conformidade com as
especificações máximas permitidas (3500kV). Além disso, suas dimensões e padrão
de montagem são compatíveis com chassis em escala 1:10, garantindo conformidade
com o regulamento~\cite{F1TENTHRules2023}. Esta restrição impõe um teto de
desempenho que:

\begin{itemize}
      \item Força as equipes a otimizarem seus algoritmos de controle para extrair o máximo
            da plataforma
      \item Minimiza as vantagens derivadas puramente de potência mecânica superior
      \item Cria condições para avaliar a eficiência dos sistemas de navegação e controle
\end{itemize}

\subsection{Sistemas de sensoriamento e computação}

O sistema de sensoriamento incorpora limitações estratégicas que simulam
restrições de orçamento e espaço em veículos reais. O LiDAR principal deve ser
equivalente ou inferior ao Hokuyo UTM-30LX, com alcance máximo de 30 metros,
frequência de varredura de 40 Hz e resolução angular de 0,25 graus
~\cite{GopalarathnamF1TENTH2022}. O LiDAR YDLIDAR X4PRO selecionado para o
projeto atende adequadamente aos requisitos regulamentares, apresentando
especificações compatíveis com os limites estabelecidos: alcance de 0,12--10m
(inferior ao limite de 30m), frequência de varredura ajustável de 6--12Hz e
resolução angular entre 0,43° e 0,86° dependendo da frequência utilizada. Sua
conformidade com a classificação Class I de segurança laser também é um aspecto
importante para atender às normas de segurança da competição
~\cite{F1TENTHRules2023}. Estas especificações restringem a quantidade e
qualidade dos dados disponíveis para percepção ambiental, exigindo:

\begin{itemize}
      \item Algoritmos eficientes de processamento de nuvem de pontos
      \item Técnicas robustas de filtragem para lidar com ruído e reflexões no ambiente de
            competição
      \item Estratégias de fusão sensorial para compensar as limitações do LiDAR
\end{itemize}

Quanto ao sistema computacional, as regulamentações permitem flexibilidade,
exigindo apenas que o processamento seja totalmente realizado a bordo e que o
hardware se adeque fisicamente ao veículo~\cite{F1TENTHRules2023}. A escolha da
Raspberry Pi 4B como plataforma de processamento satisfaz plenamente estes
requisitos, oferecendo processamento suficiente para algoritmos de navegação e
controle em um formato compacto que se adequa facilmente ao chassi do veículo.
Sua ampla compatibilidade com frameworks de software como ROS2 facilita a
implementação dos algoritmos necessários para a competição. Esta flexibilidade
reconhece a importância do poder computacional para algoritmos avançados de
percepção e controle, permitindo o uso de plataformas como Nvidia Jetson, Intel
NUC ou Raspberry Pi, conforme as necessidades específicas de processamento da
solução implementada.

O controlador eletrônico Makerbase VESC MINI 6.7 representa uma escolha
estratégica para o projeto, visto que suas capacidades avançadas de controle de
motor BLDC e PMSM, incluindo recursos de posicionamento preciso e controle de
velocidade/corrente, permitem a implementação de algoritmos sofisticados de
controle. Suas interfaces de comunicação múltiplas (USB, CAN, UART) facilitam a
integração com o sistema ROS2 executado na Raspberry Pi, criando um sistema de
controle robusto e responsivo que atende aos requisitos de processamento a
bordo exigidos pelo regulamento~\cite{F1TENTHRules2023}.

\subsection{Requisitos de segurança obrigatórios}

Os requisitos de segurança não são meras formalidades, mas elementos cruciais
que simulam preocupações de veículos autônomos em escala real. Todo veículo
deve incorporar:

\begin{itemize}
      \item Para-choques frontais e traseiros funcionais
      \item Área mínima de detecção por LiDAR de 12×12 centímetros entre 10--30 centímetros do solo
      \item Sistema de parada de emergência acionável remotamente
      \item Algoritmos de detecção e prevenção de colisões demonstráveis
\end{itemize}

A configuração do sistema VESC MINI 6.7 com capacidade de frenagem regenerativa
e limites de corrente programáveis contribui significativamente para atender
aos requisitos de segurança da competição, permitindo implementar um sistema
eficaz de parada de emergência. Sua capacidade de resposta rápida (com
limitação de corrente e proteções integradas) permite implementar algoritmos
confiáveis de detecção e prevenção de colisões, conforme exigido pelo
regulamento~\cite{MoralesF1TENTH2021}.

Estes requisitos são testados rigorosamente antes da competição, e falhas em
qualquer um deles resultam em desqualificação~\cite{MoralesF1TENTH2021}.
Particularmente para as corridas head-to-head, os veículos devem passar por
testes de evitação de obstáculos estáticos e dinâmicos, simulando situações
reais de tráfego.

\section{Ambientes de Competição e suas Características Técnicas}

O ambiente de corrida é meticulosamente projetado para apresentar desafios
específicos aos sistemas autônomos. A pista, construída em superfície plana e
reflexiva com bordas formadas por dutos de ar de 33 cm de diâmetro, cria
condições particularmente desafiadoras para sistemas de percepção
~\cite{OKelly2020F1TENTH}:

\begin{itemize}
      \item Os feixes do LiDAR frequentemente refletem no solo reflexivo, gerando ruídos e
            falsas leituras
      \item As câmeras de profundidade enfrentam dificuldades com a detecção adequada dos
            limites da pista
      \item As bordas curvas dos dutos de ar criam padrões de reflexão não-lineares e
            complexos
\end{itemize}

O LiDAR YDLIDAR X4PRO apresenta características que podem ajudar a lidar com
essas condições desafiadoras da pista, incluindo sua forte resistência a
interferências de luz ambiente (tolerância de até 40000 Lux) e capacidade de
medição estável em superfícies com diferentes refletividades. A resolução
angular configurável permite ajustes para obter o equilíbrio ideal entre
densidade de pontos e frequência de atualizações, fundamentais para a navegação
precisa nas condições específicas da competição~\cite{OKelly2020F1TENTH}.

A largura mínima de 90 centímetros (equivalente a três larguras de carro) é
especificamente calculada para permitir manobras de ultrapassagem durante as
corridas head-to-head, exigindo sistemas sofisticados de planejamento de
trajetória. O circuito combina intencionalmente elementos diversos como:

\begin{itemize}
      \item Pontos de estrangulamento que forçam decisões estratégicas de posicionamento
      \item Curvas fechadas que testam os limites dos algoritmos de controle
      \item Retas que permitem maximização de velocidade e avaliação de sistemas de
            controle longitudinal
      \item Curvas extra-largas que abrem possibilidades para múltiplas linhas de corrida
            otimizadas
\end{itemize}

A combinação do motor EZRUN 3665SD-3200KV com o controlador VESC MINI 6.7
oferece uma resposta dinâmica adequada para enfrentar estes diversos elementos
da pista. O motor fornece torque suficiente para aceleração rápida nas retas,
enquanto o controle preciso oferecido pelo VESC permite ajustes finos nas
curvas e manobras de ultrapassagem. A capacidade de programação avançada do
VESC possibilita a implementação de diferentes perfis de controle otimizados
para cada seção do circuito~\cite{BalachandranF1TENTH2020}.

Esta diversidade de elementos não é aleatória, mas cuidadosamente planejada
para testar a versatilidade dos algoritmos implementados sob diferentes
condições, simulando os variados cenários que veículos autônomos reais
enfrentam~\cite{BalachandranF1TENTH2020}.

\section{Estrutura e Funcionamento da Competição}

\subsection{Inspeção técnica e verificação de conformidade}

A competição inicia com uma rigorosa fase de inspeção técnica, onde cada
veículo é minuciosamente avaliado quanto à conformidade com as especificações
técnicas~\cite{F1TENTHRules2023}. Esta etapa inclui:

\begin{itemize}
      \item Verificação dimensional do chassi e componentes
      \item Teste dos sistemas de parada de emergência sob diferentes condições
      \item Validação das especificações do LiDAR e outros sensores
      \item Confirmação de que todo processamento ocorre a bordo
\end{itemize}

Os componentes selecionados para o projeto – motor EZRUN 3665SD-3200KV, VESC
MINI 6.7, LiDAR YDLIDAR X4PRO e Raspberry Pi 4B – foram escolhidos considerando
sua conformidade com os requisitos da inspeção técnica. As especificações
detalhadas de cada componente facilitam a documentação necessária para este
processo, enquanto as capacidades de diagnóstico do sistema VESC e a interface
de software ROS2 permitem demonstrar claramente o funcionamento dos sistemas
críticos durante a avaliação~\cite{F1TENTHRules2023}.

Veículos que não atendem integralmente às especificações são desclassificados
ou reclassificados para a Classe Aberta, dependendo da natureza da
não-conformidade.

\subsection{Time Trial: avaliação de desempenho individual}

O Time Trial constitui a segunda fase da competição, estruturada em duas
baterias de 5 minutos cada. Nesta modalidade, cada veículo compete
individualmente na pista, buscando:

\begin{itemize}
      \item Completar o maior número possível de voltas consecutivas sem interrupções
      \item Registrar a volta mais rápida possível durante o período alocado
\end{itemize}

O sistema de pontuação é deliberadamente projetado para equilibrar velocidade e
consistência, combinando o tempo da volta mais rápida com o número de voltas
consecutivas completadas sem intervenção humana~\cite{ShinF1TENTH2020}. Esta
abordagem força as equipes a desenvolverem soluções que balanceiem:

\begin{itemize}
      \item Algoritmos agressivos de otimização de trajetória para velocidade máxima
      \item Sistemas robustos de controle que garantam consistência entre múltiplas voltas
      \item Estratégias de gerenciamento de risco que minimizem a probabilidade de falhas
\end{itemize}

A integração do VESC MINI 6.7 com ROS2 conforme detalhado na documentação do
projeto proporciona um sistema de controle preciso e adaptável, crucial para o
desempenho consistente na fase de Time Trial. A capacidade de ajustar
parâmetros de controle PID e perfis de velocidade permite otimizar o veículo
tanto para voltas rápidas quanto para desempenho consistente ao longo de
múltiplas voltas, alinhando-se perfeitamente com os critérios de pontuação
desta fase~\cite{ShinF1TENTH2020}.

\subsection{Corridas Head-to-Head: o desafio máximo da autonomia}

A fase final e mais complexa é a Corrida Head-to-Head, onde dois veículos
competem simultaneamente na pista. Esta modalidade representa o ápice dos
desafios de autonomia, exigindo:

\begin{itemize}
      \item Sistemas robustos de detecção e rastreamento de outros veículos em movimento
      \item Algoritmos adaptativos de planejamento de trajetória que considerem obstáculos
            dinâmicos
      \item Estratégias sofisticadas de ultrapassagem que calculem riscos e oportunidades
      \item Capacidade de ajuste dinâmico da linha de corrida baseado na posição do
            oponente
\end{itemize}

O LiDAR YDLIDAR X4PRO, com sua capacidade de varredura de 360° e alta
frequência de amostragem (5kHz), fornece dados cruciais para a detecção de
outros veículos durante as corridas head-to-head. A resolução angular permite
identificar com precisão a posição do oponente, enquanto a Raspberry Pi 4B
oferece o poder de processamento necessário para executar algoritmos complexos
de planejamento de trajetória e tomada de decisões em tempo real, fundamentais
para implementar estratégias eficazes de ultrapassagem
~\cite{GopalarathnamF1TENTH2022}.

As corridas são organizadas em formato de eliminatórias, com os veículos mais
rápidos das fases de Time Trial recebendo posições favoráveis no grid
~\cite{GopalarathnamF1TENTH2022}. Uma característica crucial desta fase é a
necessidade de combinar velocidade máxima com tomada de decisão defensiva e
ofensiva, similar aos desafios enfrentados por veículos autônomos em ambientes
urbanos densos.

\section{Implicações das Regulamentações para o Desenvolvimento Customizado}

As regulamentações da F1TENTH impactam diretamente as decisões de projeto no
desenvolvimento de veículos customizados. Considerando o desenvolvimento de uma
versão personalizada que atenda às regulações, é necessário:

\begin{enumerate}
      \item Equilibrar o poder computacional com restrições de tamanho e energia, optando
            por plataformas que ofereçam desempenho otimizado para algoritmos específicos
            implementados. A Raspberry Pi 4B representa este equilíbrio, oferecendo
            processamento suficiente em um formato compacto e com consumo energético
            moderado compatível com as baterias típicas de veículos RC em escala 1:10.

      \item Desenvolver estratégias de fusão sensorial que compensem as limitações do LiDAR
            principal, potencialmente incorporando dados de IMUs, encoders e câmeras dentro
            dos limites regulatórios~\cite{YangF1TENTH2022}. A integração entre o LiDAR
            YDLIDAR X4PRO e os encoders disponíveis através do VESC MINI 6.7 possibilita
            implementar fusão sensorial eficaz para melhorar a percepção e navegação.

      \item Implementar sistemas de controle adaptáveis que possam alternar entre modos de
            operação otimizados para diferentes segmentos da pista e situações
            competitivas. O VESC MINI 6.7, com sua capacidade de programação avançada e
            ajuste de parâmetros em tempo real via comunicação com a Raspberry Pi, permite
            a implementação destes sistemas de controle adaptáveis.

      \item Criar algoritmos de planejamento que incorporem tanto conhecimento a priori do
            circuito quanto capacidade de adaptação a condições dinâmicas durante as
            corridas head-to-head~\cite{MoralesF1TENTH2021}

      \item Projetar sistemas de segurança que não apenas atendam aos requisitos mínimos,
            mas que também sejam suficientemente robustos para permitir operação agressiva
            próxima aos limites físicos do veículo. As proteções integradas do VESC MINI
            6.7 (sobrecorrente, sobretensão, temperatura) fornecem uma base sólida para
            implementar estes sistemas de segurança.
\end{enumerate}

As regulamentações, embora restritivas, criam um ambiente de inovação focada
onde soluções algorítmicas sofisticadas ganham precedência sobre vantagens de
hardware, refletindo adequadamente os desafios enfrentados no desenvolvimento
de veículos autônomos em escala real~\cite{OKelly2020F1TENTH}. Os componentes
selecionados para este projeto demonstram um equilíbrio adequado entre
desempenho e conformidade regulatória, proporcionando uma plataforma viável
para competir efetivamente na categoria F1TENTH.
\chapter{Considerações finais}

A primeira parte deste trabalho de conclusão de curso explorou os fundamentos
teóricos, ferramentas e regulamentações que constituem o ecossistema da
competição F1TENTH. Através desta análise, evidencia-se que a plataforma
F1TENTH representa não apenas um ambiente competitivo, mas um laboratório
educacional que sintetiza os desafios centrais da mobilidade autônoma
contemporânea.

\begin{itemize}
    \item \textbf{Integração técnica e prática:} A implementação de sistemas ROS2 em plataformas com restrições computacionais reflete fielmente os compromissos necessários no desenvolvimento de veículos autônomos comerciais, onde otimização de recursos e desempenho determinístico são requisitos inegociáveis.

    \item \textbf{Complexidade progressiva:} A estrutura da competição, evoluindo de desafios de tempo individual para confrontos diretos, cria um percurso pedagógico que gradualmente introduz as complexidades da tomada de decisão em ambientes dinâmicos.

    \item \textbf{Robustez algorítmica:} O foco em equalizar vantagens mecânicas direciona o desenvolvimento para soluções algorítmicas sofisticadas, priorizando eficiência computacional e resiliência a condições adversas.

    \item \textbf{Fusão sensorial e percepção:} As limitações deliberadas nos sistemas de sensoriamento estimulam abordagens inovadoras para fusão de dados, filtragem de ruído e extração de características em tempo real.

    \item \textbf{Prototipagem acelerada:} A escala reduzida permite ciclos rápidos de desenvolvimento e experimentação que seriam proibitivamente custosos e arriscados em veículos de tamanho real.
\end{itemize}

As restrições impostas pelas regulamentações técnicas, longe de limitarem a
inovação, criam um ambiente controlado onde soluções algorítmicas criativas
ganham precedência sobre vantagens de hardware, espelhando os compromissos
enfrentados na indústria automotiva. Esta abordagem estabelece um microcosmo
ideal para a aplicação prática dos conhecimentos teóricos de robótica móvel,
sistemas embarcados e inteligência artificial em um contexto tangível e
mensurável.

A próxima etapa deste trabalho se concentrará na implementação prática destes
conceitos, aplicando o conhecimento teórico adquirido ao desenvolvimento de um
veículo F1TENTH funcional, com ênfase particular na navegação autônoma robusta
e na operação próxima aos limites dinâmicos do sistema.

% PARTE
\part{Proposta}
\chapter{Sistema Proposto}\label{cap:proposta}

Esse trabalho propõe um sistema de... 


\section{Primeira Parte do Sistema Proposto}

\lipsum[67]

\section{Considerações Finais}

\lipsum[68]


% PARTE
\part{Parte Final}
\chapter{Resultados e Discussão}\label{cap:resultados}

\lipsum[73]

\section{Base de Dados}

\lipsum[72]

\section{Considerações Finais}

\lipsum[74]
\chapter*{Conclusões e Trabalhos Futuros}\label{cap:conclusao}
\addcontentsline{toc}{chapter}{Conclusão e Trabalhos Futuros}

\lipsum[81]

\section*{Conclusões}

\lipsum[82-84]

\section*{Trabalhos Futuros}

\lipsum[85] 

% ----------------------------------------------------------
% ELEMENTOS PÓS-TEXTUAIS (Referências, Glossário, Apêndices)
% ----------------------------------------------------------
\postextual

% Referências bibliográficas
% \makeatletter
% \let\bibstyle\@gobble
% \makeatother
\bibliography{bibliografia}

% Glossário (Consulte o manual)
%\glossary

% Apêndices
%% ----------------------------------------------------------
% Apêndices
% ----------------------------------------------------------

% ---
% Inicia os apêndices
% ---
\begin{apendicesenv}

% Imprime uma página indicando o início dos apêndices
\partapendices

% ----------------------------------------------------------
\chapter{Primeiro Apêncice}
% ----------------------------------------------------------

\lipsum[50] % Texto qualquer. REMOVER!!

% ----------------------------------------------------------
\chapter{Segundo apêndice com título tão grande quanto se queira porque ele já faz a quebra de linha da coisa toda}
% ----------------------------------------------------------
\lipsum[51-53] % Texto qualquer. REMOVER!!

\end{apendicesenv}
% ---

% Anexos
%% ----------------------------------------------------------
% Apêndices
% ----------------------------------------------------------

% ---
% Inicia os anexos
% ---
\begin{anexosenv}

% Imprime uma página indicando o início dos anexos
\partanexos

% ---
\chapter{Nome do Primeiro Anexo}
% ---
\lipsum[30] % Texto qualquer. REMOVER!!

% ---
\chapter{Nome de Outro Anexo}
% ---

\lipsum[32] % Texto qualquer. REMOVER!!

\end{anexosenv}

% Índice remissivo (Consultar manual)
%\phantompart
%\printindex

\end{document}
